%!TEX TS-program = xelatex
%!TEX encoding = UTF-8 Unicode
% ==============================================================================
% İSTANBUL ÜNİVERSİTESİ SOSYAL BİLİMLER ENSTİTÜSÜ TEZ ŞABLONU
% Bu şablon İ.Ü. SBE Tez Hazırlama Yönergesi esas alınarak hazırlanmıştır.
% Derleme Motoru: XeLaTeX (veya LuaLaTeX) + Biber (kaynakça için)
% ==============================================================================
\documentclass[12pt,a4paper,oneside]{iusosbil}

% İhtiyaç duyulabilecek ek paketler
\usepackage{lipsum}      % Test amaçlı sahte metin üretici (tez tesliminde kaldırınız)
\usepackage{egyphitt}    % Hititçe, Mısıroloji ve Çivi Yazısı desteği

\begin{document}
\justifying              % Madde 4/b: Satırları iki yana yasla

% ==============================================================================
% A. ÖN KISIM (Madde 8/A)
% ==============================================================================

% 1. İç Kapak (Madde 9 & Örnek 1 - Sayfa numarası görünmez, sayfa i kabul edilir)
%!TEX root = main.tex
% ==============================================================================
% İSTANBUL ÜNİVERSİTESİ SOSYAL BİLİMLER ENSTİTÜSÜ
% İÇ KAPAK (Madde 9 ve Ek Madde Örnek 1)
% ==============================================================================
\begin{titlepage}
\begin{center}
    \bfseries
    \fontsize{14pt}{18pt}\selectfont

    T.C.\\
    İSTANBUL ÜNİVERSİTESİ\\
    SOSYAL BİLİMLER ENSTİTÜSÜ\\
    ANABİLİM DALI ADI\\
    BİLİM DALI ADI

    \vspace{2cm}

    \fontsize{14pt}{18pt}\selectfont
    YÜKSEK LİSANS TEZİ / DOKTORA TEZİ

    \vspace{2cm}

    % Başlık: 18 punto (Örnek 1)
    \fontsize{18pt}{22pt}\selectfont
    TEZİN BAŞLIĞI BURAYA GELECEK

    \vspace{2.5cm}

    \fontsize{14pt}{18pt}\selectfont
    Adı SOYADI\\
    Öğrenci Numarası: 123456789

    \vfill

    \fontsize{14pt}{18pt}\selectfont
    Tez Danışmanı\\
    Prof. Dr. Adı SOYADI

    \vspace{1.5cm}

    \fontsize{14pt}{18pt}\selectfont
    İstanbul 20XX

\end{center}
\end{titlepage}


%!TEX root = main.tex
% ==============================================================================
% İSTANBUL ÜNİVERSİTESİ SOSYAL BİLİMLER ENSTİTÜSÜ
% TEZ ONAY SAYFASI (Madde 10)
% Madde 3/a uyarınca bu sayfada sayfa numarası yer almaz.
% ==============================================================================
\thispagestyle{empty}
\begin{center}
    \bfseries
    \fontsize{14pt}{18pt}\selectfont
    TEZ ONAY SAYFASI
\end{center}

\vspace{1.5cm}

\noindent
İstanbul Üniversitesi Sosyal Bilimler Enstitüsü \textbf{..................................} Anabilim Dalı \textbf{..................................} Bilim Dalı'nda \textbf{......................} numaralı \textbf{..................................}'nın hazırladığı \textbf{“................................................”} konulu \textbf{YÜKSEK LİSANS / DOKTORA TEZİ}, jürimiz tarafından \textbf{..... / ..... / 20XX} tarihinde yapılan tez savunma sınavında başarılı bulunarak kabul edilmiştir.

\vspace{2.5cm}

\begin{center}
\textbf{TEZ JÜRİSİ ÜYELERİ}
\end{center}

\vspace{1.5cm}

\noindent
\begin{tabular*}{\textwidth}{@{\extracolsep{\fill}}llr}
    \textbf{Unvanı, Adı ve Soyadı} & \textbf{Üniversitesi} & \textbf{İmza} \\[1.3cm]
    1. Tez Danışmanı: ............................ & ............................ & ................ \\[1.3cm]
    2. Jüri Üyesi: ..................................... & ............................ & ................ \\[1.3cm]
    3. Jüri Üyesi: ..................................... & ............................ & ................ \\[1.3cm]
    4. Jüri Üyesi: ..................................... & ............................ & ................ \\[1.3cm]
    5. Jüri Üyesi: ..................................... & ............................ & ................ \\
\end{tabular*}

\clearpage
 % 2. Tez Onay Sayfası (Madde 10 - Sayfa numarası görünmez, sayfa ii kabul edilir)
% 3. Öz (iii), Abstract (iv), Önsöz, İçindekiler, Tablo/Şekil/Kısaltma Listeleri
% Madde 3/a uyarınca numaralar sağ alt köşede (iii) ile başlar.
%!TEX root = main.tex
% ==============================================================================
% İSTANBUL ÜNİVERSİTESİ SOSYAL BİLİMLER ENSTİTÜSÜ
% ÖN KISIM (Madde 8/A)
% ==============================================================================

% Madde 3/a: Numaralama iç kapaktan başlar ancak ilk gösterilen sayfa Öz sayfasıdır ve (iii)'tür.
\baslatOnKisim

% ------------------------------------------------------------------------------
% 1. ÖZ (TÜRKÇE) (Madde 11)
% Sayfa başında tezin başlığı ve yazar adı yer almalıdır (250 kelimeyi geçmemelidir).
% İtalik yazı, simge, formül içermemelidir (Madde 5/c).
% ------------------------------------------------------------------------------
\chapter*{ÖZ}
\addcontentsline{toc}{chapter}{ÖZ}

\begin{center}
    {\bfseries\fontsize{14pt}{16pt}\selectfont TEZİN BAŞLIĞI BURAYA GELECEK}\\[0.4cm]
    {\bfseries Adı SOYADI}
\end{center}

\vspace{0.8cm}

Bu çalışmada, İstanbul Üniversitesi Sosyal Bilimler Enstitüsü tez hazırlama yönergesine uygun olarak belirlenen araştırma problemi, yöntemi ve elde edilen temel bulgular özetlenmektedir. Araştırmada incelenen sorunsal, kullanılan kavramsal çerçeve, yöntem ve teknikler ile varılan sonuçlar bu kısımda özlü ve açık bir dille ortaya konulmuştur. Öz metni 250 kelimeyi aşmayacak şekilde yapılandırılmıştır.

\vspace{1cm}
\noindent\textbf{Anahtar Kelimeler:} Sosyal Bilimler, Tez Yazımı, İstanbul Üniversitesi, LaTeX Şablonu, Akademik Standartlar.

\clearpage

% ------------------------------------------------------------------------------
% 2. ABSTRACT (İNGİLİZCE) (Madde 11)
% İngilizce başlık ve yazar adı ile başlar (250 kelimeyi geçmemelidir).
% ------------------------------------------------------------------------------
\begin{otherlanguage}{english}
\chapter*{ABSTRACT}
\addcontentsline{toc}{chapter}{ABSTRACT}

\begin{center}
    {\bfseries\fontsize{14pt}{16pt}\selectfont THESIS TITLE IN ENGLISH}\\[0.4cm]
    {\bfseries Author NAME SURNAME}
\end{center}

\vspace{0.8cm}

In this thesis, the research questions, methodology, and primary findings are summarized in accordance with the Istanbul University Institute of Social Sciences thesis preparation guidelines. The conceptual framework, qualitative/quantitative methods, and conclusive findings are presented concisely without exceeding 250 words.

\vspace{1cm}
\noindent\textbf{Keywords:} Social Sciences, Thesis Guidelines, Istanbul University, LaTeX Template, Academic Standards.
\end{otherlanguage}

\clearpage

% ------------------------------------------------------------------------------
% 3. ÖNSÖZ (Madde 12)
% Tezin amacı, önemi, karşılaşılan güçlükler ve teşekkür içerir. En fazla 2 sayfadır.
% ------------------------------------------------------------------------------
\chapter*{ÖNSÖZ}
\addcontentsline{toc}{chapter}{ÖNSÖZ}

Bu tez çalışması, İstanbul Üniversitesi Sosyal Bilimler Enstitüsü bünyesinde yürütülmüş olup alan yazına katkı sağlamayı amaçlamaktadır.

Çalışmanın her aşamasında değerli bilgi, öneri ve eleştirileriyle yol gösteren, desteğini esirgemeyen saygıdeğer danışmanım Sayın Prof. Dr. Adı SOYADI'na en içten teşekkürlerimi sunarım.

Ayrıca, araştırmam süresince maddi ve manevi desteklerini daima yanımda hissettiğim aileme ve çalışma arkadaşlarıma teşekkür ederim.

\vspace{1.5cm}
\begin{flushright}
    Adı SOYADI\\
    İstanbul, 20XX
\end{flushright}

\clearpage

% ------------------------------------------------------------------------------
% 4. İÇİNDEKİLER (Madde 13)
% Öz'den başlayarak tezin tüm kısımlarını sayfa numaralarıyla listeler.
% ------------------------------------------------------------------------------
\tableofcontents

\clearpage

% ------------------------------------------------------------------------------
% 5. TABLOLAR VE ŞEKİLLER LİSTELERİ (Madde 14)
% ------------------------------------------------------------------------------
\addcontentsline{toc}{chapter}{\listfigurename}
\listoffigures

\clearpage

\addcontentsline{toc}{chapter}{\listtablename}
\listoftables

\clearpage

% ------------------------------------------------------------------------------
% 6. KISALTMALAR LİSTESİ (Madde 15)
% Alfabetik sırayla alt alta dizilir, karşısında açık halleri yer alır.
% ------------------------------------------------------------------------------
\chapter*{KISALTMALAR}
\addcontentsline{toc}{chapter}{KISALTMALAR}

\begin{description}[leftmargin=3.5cm, style=nextline]
    \item[a.e.] Aynı eser
    \item[a.g.e.] Adı geçen eser
    \item[a.y.] Aynı yer
    \item[bs.] Basım
    \item[bsk.] Baskı
    \item[C.] Cilt
    \item[Çev.] Çeviren
    \item[Ed.] Editör
    \item[İ.Ü.] İstanbul Üniversitesi
    \item[No.] Numara
    \item[s.] Sayfa
    \item[SBE] Sosyal Bilimler Enstitüsü
    \item[T.C.] Türkiye Cumhuriyeti
    \item[t.y.] Tarih yok
    \item[v.d.] Ve diğerleri
    \item[y.y.] Yayım yeri yok / Yayımcı yok
\end{description}



% ==============================================================================
% B. METİN KISMI (Madde 8/B)
% Giriş'ten itibaren Arap rakamlarıyla (1, 2, 3...) numaralandırma başlar (Madde 16/a).
% ==============================================================================
\baslatMetin
%!TEX root = main.tex
% ==============================================================================
% İSTANBUL ÜNİVERSİTESİ SOSYAL BİLİMLER ENSTİTÜSÜ
% GİRİŞ BÖLÜMÜ (Madde 16/a, b ve Madde 17)
% Giriş, metin kısmının ilk öğesidir; Arap rakamlarıyla numaralandırma buradan başlar.
% Giriş bölümü numaralanmaz (Madde 17).
% ==============================================================================
\chapter*{GİRİŞ}
\addcontentsline{toc}{chapter}{GİRİŞ}

Giriş kısmında, önsözde belirtilen hususlar tekrar edilmemek üzere, çalışmada çözümlenmesi amaçlanan bilimsel sorun/sorunsal etraflıca tanımlanır. Bu doğrultuda araştırmada yararlanılan kavramsal çerçeve, kuramsal dayanaklar, yöntem, teknik ve varsa kullanılan paradigmalar ayrıntılı bir biçimde açıklanır (Madde 16/b).

Tez metninde satırlar iki yana yaslanmış ve 1,5 satır aralığı ile yazılmıştır (Madde 4/b). Paragraflar arasında gereksiz dikey boşluk bırakılmamış, standart 1 cm paragraf girintisi uygulanmıştır. 

Araştırma konusuna ilişkin genel değerlendirmeler ve tezin kurgusu Giriş kısmında ortaya konulduktan sonra, tezi sonuca ulaştıracak temel bölümlere geçilmektedir.

 % Giriş Bölümü (Madde 16/a, b) - Numarasızdır
%!TEX root = main.tex
% ==============================================================================
% İSTANBUL ÜNİVERSİTESİ SOSYAL BİLİMLER ENSTİTÜSÜ
% BÖLÜM 1 (Madde 16/c, d ve Madde 17/b - İkinci Sistem)
% ==============================================================================
\chapter{KURAMSAL ÇERÇEVE VE KAVRAMSAL YAKLAŞIMLAR}

\section{Temel Kavramlar ve Yöntemsel İlkeler}

Sosyal bilim araştırmalarında kullanılan terimlerde ve üslupta birlik sağlanması yönergenin temel ilkelerindendir. Örneğin, metinde aynı anlama gelen terimlerin (örneğin ‘yasa’ / ‘kanun’ ya da ‘kuram’ / ‘teori’) hangisi seçildiyse o terimde tutarlılık gösterilmelidir.

Başka kaynaklara yapılan atıflar, araştırmacının tercihine göre metin içi yazar-yıl formatında \parencite[11]{sartori_demokrasi} veya sayfa altı referans dipnotu şeklinde verilebilir.\footnote{Giovanni Sartori, \textbf{Demokrasi Kuramı}, Çev. Deniz Baykal, Ankara, Siyasi İlimler Türk Derneği Yayınları, t.y., s. 11.} Yönerge Madde 19/a uyarınca, dipnot numaraları her bölüm başında yeniden 1'den başlar ve tek satır aralıklı, 10 punto ile dizilir.\footnote{Aynı esere yapılan sonraki atıflarda kısaltmalar (a.g.e., a.e. veya op. cit.) tutarlı olarak kullanılır: Sartori, \textbf{a.g.e.}, s. 22.}

\subsection{Doğrudan Alıntılar ve Blok Aktarma Kuralları}

Bir cümleyi aşmayan kısa alıntılar metin içinde çift tırnak (`"..."') ile gösterilir (Madde 18/a). Daha uzun aktarmalarda ise alıntı ile metin arasında belirgin bir ayrım yapılmalıdır (Madde 18/e ve Madde 4/b):

\begin{uzunalinti}
Kaynaklardan aynen yapılan alıntılar, bir cümleyi aşmayacak uzunluktaysa, tez metni içinde verilirler. Daha uzun alıntılarda, alıntının ilk ve son satırları ile tez metni arasında en az çift aralıklı boşluk bırakılarak, alıntı ile metnin ayırt edilmesi sağlanır. Bu tür alıntılar için, tez metninde kullanılan puntodan daha küçük punto kullanılır ve satırlar tek aralıklı yazılır (Madde 4/b, 18/e).
\end{uzunalinti}

\section{Özel Karakterler ve Transliterasyon}
Tezde yabancı alfabeler, hiyeroglif veya çivi yazısı transliterasyonları kullanılması durumunda uygun yazı tipleri devreye girer:
\begin{center}

\vspace{0.8em}

\translit{1-jr m ntw.f ḫr.tw sn.wy n wꜥ m(jw)t n wꜥ jt(j) jnpw rn pꜣ ꜥꜣ jw bꜣtꜣw rn pꜣ šrjw ḫr jr jnpw sw ẖry pr ẖry ḥmt {}
  2-jw pꜣy.f sn šrjw m-dj.f mj sḫrw n šrjw jw m ntw.f jrr.tw n.f ḥbsw jw.f ḥr šmyt m-sꜣ nꜣy.f jꜣwt r sḫtw \\{}
  3-[j]w m ntw.f jrr.tw skꜣ m ntw.f ꜥwꜣyw ḫr m ntw.f jrrt.n.f wpwt nb nty m sḫtw js.tw jw pꜣy.f sn šrjw} \\{}
\end{center}

\vspace{0.8em}

\begin{htranslit}
1-[LUGAL.GAL Ta-ba-ar]-⸢na⸣ m Ḫa-at-tu-ši-li LUGAL.GAL 
2-[LUGAL KUR URU Ḫa-at-]ti LÚ URU Ku-uš-šar KUR URUḪa-at-⸢ti⸣ 
3-[LUGAL e-ez-zi-a-at] ŠA fTa-wa-an-na-an-na DUMU ŠEŠ-ŠU  
4-[I-NA URU Š]a-na-wi-it-ta pa-it ša-an na-at-ta 
5-[ḫar-ni-ik]-⸢t⸣a na ut-ni-eš še-it ḫar-ni-ik-ta 
6-[nu-kán ERÍN.MEŠ] 2 AŠ-RA a-ša-an-du-la-ni da-<la>-⸢aḫ-ḫu⸣-un 
7-[ku-u]-e a-ša-u-wa-ar e-eš-[t]a 
8-[na-at A-⸢NA⸣] ERÍN.MEŠ a-ša-an-du-li pí-aḫ-ḫu-⸢un⸣ 
\tabletrule
9-[EGIR-an-da]-ma I-NA URU Za-la-pa pa-a[..............]-un |] \\
10-[na-an]⸢ḫa⸣r-ni-in-ku-un nu-uš-ši DMEŠ-ŠU ša-ra-a da-aḫ-ḫu-un|] \\
11-[Ù 3 GIŠ GIGIR.MEŠ MA-AD-NA-NU A-NA D.UTU URUTÚL-na pí-aḫ-ḫu-un|]
\tabletrule
12-[1 GUD KÙ.BABBAR 1 GEŠPU KÙ.BABBAR I-NA É D IM pí-aḫ-ḫu-un |] \\
13-[a-aš-ir-ma]-kán ku-i-e-eš na-aš I-NA É D Me-ez-zu-ul-la |] \\
14-[pí-da-aḫ-ḫu[...............................................]-un |]
  
\end{htranslit}

Bu sayede eskiçağ dilleri ve filolojik metinler özgün işaretleriyle bozulmadan aktarılabilmektedir.

 % Tezin Bölümleri (Madde 16/c, d - Giriş hariç en az üç bölümden oluşur)
%!TEX root = main.tex
% ==============================================================================
% İSTANBUL ÜNİVERSİTESİ SOSYAL BİLİMLER ENSTİTÜSÜ
% BÖLÜM 2 (Madde 16/c, d ve Madde 17/b)
% ==============================================================================
\chapter{VERİ ANALİZİ VE BULGULAR}

\section{Araştırma Bulgularının Sınıflandırılması}

Bu bölümde toplanan ampirik veriler ve kuramsal karşılaştırmalar yer almaktadır.\footnote{Yeni bölüme geçildiğinde dipnot numaralandırması Madde 19/a gereğince tekrar 1'den başlar.} Sosyolojik sistem analizlerinde sıklıkla başvurulan yaklaşımlardan biri Parsons tarafından geliştirilmiştir \parencite{parsons_social_system}.

\begin{table}[htbp]
    \centering
    \caption{Araştırmada İncelenen Değişkenlerin Dağılımı}
    \label{tab:degiskenler}
    \begin{tabular}{|l|c|r|}
        \hline
        \textbf{Kategori} & \textbf{Örneklem Sayısı (N)} & \textbf{Oran (\%)} \\
        \hline
        Birinci Grup & 120 & 40,0 \\
        İkinci Grup  & 150 & 50,0 \\
        Üçüncü Grup  & 30  & 10,0 \\
        \hline
        \textbf{Toplam} & \textbf{300} & \textbf{100,0} \\
        \hline
    \end{tabular}
\end{table}

Tablo \ref{tab:degiskenler} incelendiğinde, örneklem gruplarının oransal dağılımı görülmektedir. Tablo ve şekillerin numaralandırılmasında Arap veya Romen rakamları tutarlı biçimde kullanılır (Madde 3/b).

\section{Kurumsal Yapı ve Karşılaştırmalı Değerlendirme}

Türkiye'de anayasal gelişmeler ve insan hakları kuramları incelenirken yerli kaynaklara ve süreli yayınlara başvurulmuştur \parencite{tanor_osmanli, kaboglu_dusunce}.

\begin{figure}[htbp]
    \centering
    \framebox[0.75\textwidth][c]{\rule{0pt}{4cm} Şekil / Grafik Alanı}
    \caption{Araştırma Modeli ve Değişkenler Arası Etkileşim}
    \label{fig:model}
\end{figure}

Şekil \ref{fig:model}'de görüldüğü üzere modelin kavramsal boyutları birbiriyle bütüncül bir etkileşim içerisindedir. Şekil ve tablo listeleri ön kısımda tek satır aralığı ile yer almaktadır (Madde 4/b).
%!TEX root = main.tex
% ==============================================================================
% İSTANBUL ÜNİVERSİTESİ SOSYAL BİLİMLER ENSTİTÜSÜ
% BÖLÜM 3 (Madde 16/c, d ve Madde 17/b)
% Madde 16/d: Giriş kısmı hariç her tez en az üç bölümden oluşmalıdır.
% ==============================================================================
\chapter{TARTIŞMA VE ELEŞTİREL DEĞERLENDİRME}

\section{Bulguların Alanyazın Işığında Değerlendirilmesi}

Elde edilen bulgular, daha önce yapılmış kuramsal ve ampirik çalışmalarla karşılaştırmalı olarak tartışılmaktadır.\footnote{Her bölüm başında olduğu gibi dipnot sayacı 1'den başlar.} Hukuk felsefesi ve sosyolojisinde iktidar ilişkilerinin analizi için Hirsch'in değerlendirmeleri dikkate değerdir \parencite{hirsch_iktidar}.

\subsection{Alt Boyutlar ve Kuramsal Yansımalar}

Bu alt başlıkta, kuram ile uygulama arasındaki uyum veya farklılaşma noktaları irdelenmektedir.

\subsubsection{Birinci Alt Boyut: Yapısal Faktörler}

Yapısal unsurlar, kurumsal mekanizmaların işleyişi üzerinde belirleyici role sahiptir.

\subsubsection{İkinci Alt Boyut: İşlevsel Dinamikler}

İşlevsel dinamikler ise kurumsal dönüşümlerin hızını ve yönünü tayin etmektedir.

\section{Çalışmanın Sınırları ve Yöntemsel Kısıtlar}

Araştırma sürecinde ulaşılan kaynaklar ve veri toplama imkânları çerçevesinde ortaya çıkan yöntemsel kısıtlılıklar bu kısımda açıklanmıştır.




%!TEX root = main.tex
% ==============================================================================
% İSTANBUL ÜNİVERSİTESİ SOSYAL BİLİMLER ENSTİTÜSÜ
% SONUÇ BÖLÜMÜ (Madde 16/e ve Madde 17)
% Giriş gibi Sonuç bölümü de numaralandırılmaz.
% ==============================================================================
\chapter*{SONUÇ VE ÖNERİLER}
\addcontentsline{toc}{chapter}{SONUÇ VE ÖNERİLER}

Bu bölümde, girişte açıklanan tezin amacı ve hipotezlerinden hareketle tezin yöntemi, tekniği ve sınırlılıkları çerçevesinde ulaşılan çözümler, tezin bölümlerinde varılan sonuçlarla bütünleştirilerek açıklanmıştır (Madde 16/e).

Araştırma sonucunda elde edilen temel çıkarımlar şunlardır:
\begin{enumerate}
    \item Birinci temel bulgu ve kuramsal katkı.
    \item İkinci temel çıkarım ve alanyazına yönelik pratik katkılar.
    \item Yöntemsel süreçte elde edilen deneyimler ve sonraki araştırmacılara yönelik öneriler.
\end{enumerate}

Ayrıca yapılan araştırmada çözümlenemeyen veya kapsam dışı bırakılan meselelerin, gelecekte hangi tür ve konudaki araştırmalarla ele alınabileceğine ilişkin önerilere de yer verilmiştir. % Sonuç Bölümü (Madde 16/e) - Numarasızdır
%!TEX root = main.tex
% ==============================================================================
% İSTANBUL ÜNİVERSİTESİ SOSYAL BİLİMLER ENSTİTÜSÜ
% KAYNAKÇA / BİBLİYOGRAFYA (Madde 22)
% Kaynaklar yazarların soyadına göre alfabetik olarak sıralanır.
% ==============================================================================

% BibLaTeX ile otomatik kaynakça basımı (Madde 22 uyarınca alfabetik dizilir):
\printbibliography[heading=bibintoc, title={KAYNAKÇA}]

% NOT: Eğer BibLaTeX yerine elle thebibliography ortamı kullanılmak istenirse:
% \chapter*{KAYNAKÇA}
% \addcontentsline{toc}{chapter}{KAYNAKÇA}
% \begin{thebibliography}{99}
%   \bibitem{sartori} Sartori, Giovanni: \textbf{Demokrasi Kuramı}, Çev. Deniz Baykal, Ankara, Siyasi İlimler Türk Derneği Yayınları, t.y.
%   \bibitem{parsons} Parsons, Talcott: \textbf{The Social System}, London, Routledge and Kegan Paul, 1964.
%   \bibitem{tanor} Tanör, Bülent: \textbf{Osmanlı-Türk Anayasal Gelişmeleri}, İstanbul, Der, 1992.
% \end{thebibliography}

 % Kaynakça / Bibliyografya (Madde 22) % Ekler (Madde 23)
%!TEX root = main.tex
% ==============================================================================
% İSTANBUL ÜNİVERSİTESİ SOSYAL BİLİMLER ENSTİTÜSÜ
% EKLER (Madde 23)
% ==============================================================================
\chapter*{EKLER}
\addcontentsline{toc}{chapter}{EKLER}

Tez metninin ekler kısmında; metinle ilgili anket formları, mülakat soru listeleri, tablolar, arşiv belgeleri, haritalar ve belgelerin kopyaları yer alır (Madde 23). Ekler birden çok ise Arap veya Romen rakamlarıyla birbirini izleyecek biçimde numaralandırılır ve her ek için açıklayıcı bir başlık konulur.

\section*{EK 1: Anket Formu ve Soru Cetveli}
\addcontentsline{toc}{section}{EK 1: Anket Formu ve Soru Cetveli}
Araştırma kapsamında katılımcılara yöneltilen anket soruları burada yer alır.

\vspace{1cm}

\section*{EK 2: İzin ve Onay Belgeleri}
\addcontentsline{toc}{section}{EK 2: İzin ve Onay Belgeleri}
Etik kurul onayı ve araştırma iznine ilişkin resmi yazışmalar bu kısma eklenir.% Özgeçmiş (Madde 8/B ve Madde 3 - Doktora tezleri için zorunludur)
%!TEX root = main.tex
% ==============================================================================
% İSTANBUL ÜNİVERSİTESİ SOSYAL BİLİMLER ENSTİTÜSÜ
% ÖZGEÇMİŞ (Madde 8/B ve Madde 3)
% Yönerge gereğince sadece Doktora tezlerinde yer alır; metin numaralandırmasının
% son kısmını oluşturur. Yüksek Lisans tezlerinde bu bölüm dahil edilmez.
% ==============================================================================
\chapter*{ÖZGEÇMİŞ}
\addcontentsline{toc}{chapter}{ÖZGEÇMİŞ}

\noindent\textbf{Kişisel Bilgiler}
\begin{itemize}[leftmargin=2cm]
    \item[Adı Soyadı:] Adı SOYADI
    \item[Doğum Yeri ve Yılı:] İstanbul, 19XX
    \item[E-posta:] ornek@istanbul.edu.tr
\end{itemize}

\vspace{0.5cm}
\noindent\textbf{Eğitim Durumu}
\begin{itemize}[leftmargin=2cm]
    \item[Lisans:] İstanbul Üniversitesi, Edebiyat Fakültesi, 20XX
    \item[Yüksek Lisans:] İstanbul Üniversitesi, Sosyal Bilimler Enstitüsü, 20XX
    \item[Doktora:] İstanbul Üniversitesi, Sosyal Bilimler Enstitüsü, 20XX
\end{itemize}

\vspace{0.5cm}
\noindent\textbf{Mesleki Deneyim}
\begin{itemize}[leftmargin=2cm]
    \item[20XX - Günümüz:] Araştırma Görevlisi, İstanbul Üniversitesi
\end{itemize}

\vspace{0.5cm}
\noindent\textbf{Yayınlar ve Bildiriler}
\begin{itemize}[leftmargin=2cm]
    \item Soyadı, A. (20XX). "Makale Başlığı", \textit{Akademik Dergi Adı}, C. 1, S. 2, s. 10-25.
\end{itemize}

\end{document}
